% ============================================================
%  P4 -- An atomistic calibration anchor for a ballistic GLAD
%        simulator: Cu(100) adatom energetics and the limits of a
%        grid-local surface-diffusion kernel
%
%  Template: AIP Publishing (revtex4-1, aip class option)
%  Target: short methods paper / Computational Materials Science / JVST-A
%
%  FIRST DRAFT 2026-09-10 -- generated from:
%    01_GLAD_SIMULATION/P4_EAM_MD_FUTURE_WORK/{P4_LAMMPS_VALIDATION_20260909.md,
%      P4_HOP_BUDGET_CALIBRATION_RESULT_20260909.md, P4_DIFFUSION_LENGTH_CALIBRATION_20260908.md,
%      P4_D3_LONGRANGE_DIFFUSION_KERNEL_SCOPE_20260908.md}
%    01_GLAD_SIMULATION/simulation_batch/runs/P1_DIFFUSION_DECOMP_C2_KERNELFIX_20260907/C2_ANALYSIS_20260910.md
%    05_THESIS_PAPER_ASSETS/P4_EAM_MD_MANUSCRIPT_SKELETON_20260910.md
%  Numbers not recomputed. SCOPE NOTE: this is the "anchor + kernel-limits" paper
%  that is viable now; a "calibrated ballistic GLAD run" needs the D3 kernel
%  redesign first. TODO(author): figures, venue, decide whether to wait for D3.
% ============================================================
\documentclass[aip,jcp,amsmath,amssymb,reprint]{revtex4-1}
\usepackage{graphicx}
\usepackage{siunitx}
\usepackage{booktabs}

\begin{document}

\title{An atomistic calibration anchor for a ballistic GLAD simulator:
Cu(100) adatom energetics and the limits of a grid-local surface-diffusion
kernel}

\author{Oussama Ben Haj Ali}
\affiliation{Laboratoire de Photovolta\"{i}que et Mat\'{e}riaux Semi-conducteurs (LPMS),
\'{E}cole Nationale d'Ing\'{e}nieurs de Tunis (ENIT), Universit\'{e} de Tunis El Manar, Tunisia}
\author{Bruno Gallas}
\affiliation{Institut des NanoSciences de Paris (INSP), CNRS, Sorbonne Universit\'{e}, Paris, France}
\author{Ferid Chaffar Akkari}
\affiliation{Laboratoire de Photovolta\"{i}que et Mat\'{e}riaux Semi-conducteurs (LPMS),
\'{E}cole Nationale d'Ing\'{e}nieurs de Tunis (ENIT), Universit\'{e} de Tunis El Manar, Tunisia}

\date{\today}

\begin{abstract}
Ballistic glancing-angle-deposition (GLAD) simulators neglect surface diffusion;
knowing when that matters requires a physical length scale. We compute the
Cu(100) adatom hop barrier by nudged elastic band on an embedded-atom-method
potential, obtaining $E_a = \SI{0.5106}{eV}$, in agreement with published values
to \SIrange{0.02}{0.03}{eV}. Combined with the deposition rate this gives a
room-temperature adatom diffusion length before burial of
$L_D \approx \SIrange{20}{30}{nm}$ -- comparable to the inter-column spacing at
high deposition angle. We then show that the ballistic simulator's built-in
surface-diffusion kernel, a fixed three-by-three grid-cell contact search with a
spatial reach of about \SI{1.5}{nm}, cannot represent diffusion on the $L_D$
scale by any parameter choice: a geometry--kinetics decomposition finds the
kernel produces a bounded, single-parameter correction (film mass and
connectivity preserved, column tilt unchanged to within \SI{1}{\degree}) but no
route to the redistribution of mass between columns that $L_D \sim$
inter-column-spacing implies. The atomistic measurement therefore serves both as
a calibration anchor and as a specification for a redesigned, range-extended
diffusion kernel.
\end{abstract}

\maketitle

% ============================================================
\section{Introduction}
\label{sec:intro}
A purely ballistic GLAD growth model -- straight vapour trajectories, sticking on
contact, no adatom motion -- reproduces the tilted porous columnar morphology of
real films at high deposition angle, and its scope limits are well characterised.
One omission is surface diffusion. Whether neglecting it is acceptable for a
given quantity depends on how far an adatom moves before it is buried, relative
to the feature size of interest. That length is not available from the ballistic
model itself. Here we measure it atomistically for Cu, use it to bound the
regime in which the ballistic model is complete, and test whether the model's
own optional diffusion kernel can be calibrated to it.

% ============================================================
\section{Methods}
\label{sec:methods}

\subsection{Atomistic hop barrier}
The Cu(100) adatom hop barrier is computed by the nudged elastic band method on
the Mishin \textit{et al.}~\cite{mishin2001} embedded-atom-method potential, using a
standard slab geometry and a single adatom hopping between adjacent fourfold
hollow sites.

\subsection{Diffusion-length estimate}
The hop rate is $\nu = \nu_0 \exp(-E_a/k_BT)$ with $k_BT = \SI{25.69}{meV}$ at
\SI{298}{K}. The burial time is $t_\mathrm{bury} = d_\mathrm{ML}/R_\mathrm{dep}$
with $d_\mathrm{ML} = \SI{0.209}{nm}$ and $R_\mathrm{dep} = \SI{0.3}{nm/s}$. The
number of hops before burial is $N = \nu\,t_\mathrm{bury}$ and the
root-mean-square lateral diffusion length is
$L_D = a_\mathrm{nn}\sqrt{N}$ with $a_\mathrm{nn} = \SI{0.2556}{nm}$.

\subsection{Kernel probe}
The ballistic engine's diffusion kernel moves a newly deposited atom to the best
contact site within a fixed search radius; its contact search is a hardcoded
three-by-three grid-cell scan with cell size \SI{1}{nm}, giving an effective
reach of about \SI{1.5}{nm}. Its behaviour was probed by (i) a geometry--kinetics
decomposition comparing diffusion-on runs at dwell-hop cap 1 to matched
diffusion-off ballistic references, and (ii) a hop-budget sweep.

% ============================================================
\section{Results}
\label{sec:results}

\subsection{The anchor}
\label{ssec:anchor}
The computed barrier is $E_a = \SI{0.5106}{eV}$, consistent with
Boisvert and Lewis~\cite{boisvert1997} (\SI{0.49}{eV}, their own EAM/MD
Arrhenius fit) and Mehl \textit{et al.}~\cite{mehl1999} model-II
(\SI{0.487}{eV}) to \SIrange{0.02}{0.03}{eV}; Zhang \textit{et
al.}~\cite{zhang2011}'s isolated-adatom value (\SI{0.66}{eV}) is the
outlier. The resulting diffusion length is

\begin{table}[t]
\caption{Room-temperature Cu adatom diffusion length before burial.}
\label{tab:LD}
\begin{ruledtabular}
\begin{tabular}{ccc}
$\nu_0$ (Hz) & hops before burial & $L_D$ (rms) \\
\colrule
\num{1e12} & \num{1.6e3} & \SI{10}{nm} \\
\num{5e12} & \num{8.1e3} & \SI{23}{nm} \\
\num{1e13} & \num{1.6e4} & \SI{33}{nm} \\
\end{tabular}
\end{ruledtabular}
\end{table}

with a central estimate $L_D \approx \SIrange{20}{30}{nm}$. This is comparable
to the inter-column spacing at high deposition angle, i.e.\ room-temperature Cu
surface diffusion should be able to redistribute mass between columns, not merely
settle adatoms locally.

\subsection{The kernel cannot reach that scale}
\label{ssec:kernel}
The geometry--kinetics decomposition at the one height-matched deposition angle
(\ang{89}, \SI{300}{nm}) shows the diffusion kernel at dwell-hop cap 1 produces a
bounded, single-parameter correction: the largest-component fraction goes from
\num{0.997} to \num{0.956} (both connected, both percolating in both directions),
atom count rises \SI{14}{\percent} (diffusion fills in, it does not disperse),
interpenetration does not worsen, and column tilt is unchanged to within
\SI{1}{\degree}. Lower deposition angles are directionally consistent. Separately,
a hop-budget sweep finds the kernel is effectively binary: at budget $\le 2$ the
effect is mild, and above it the film over-densifies catastrophically, with no
intermediate regime reachable by parameter change. The reason is structural: the
contact search reaches about \SI{1.5}{nm}, one to two orders of magnitude below
$L_D$, so no setting lets it represent $L_D$-scale transport.

% ============================================================
\section{Discussion}
\label{sec:discussion}
The atomistic measurement does two things. As a calibration anchor it fixes the
physical diffusion length the ballistic model should be judged against, and it
bounds the regime -- low deposition angle, quantities insensitive to
between-column mass transport -- where the diffusion-free model is complete. As a
specification it shows the current kernel is the wrong tool for the
$L_D \sim \SI{20}{30}{nm}$ regime: a redesigned kernel with an expanding-ring
contact search, whose reach is set by $L_D$ rather than by a fixed grid stencil,
is required before the anchor can be used inside a calibrated ballistic run.

A separate short-timescale molecular-dynamics deposition test measures mechanical
settling of a landed atom into the nearest favourable site, on a picosecond
window; that is a different quantity from thermally activated surface diffusion,
which needs milliseconds to seconds of real time to accumulate the
\num{1e3}--\num{1e4} hops behind $L_D$. The two should not be combined into a
single calibrated mobility parameter.

% ============================================================
\section{Conclusion}
\label{sec:conclusion}
A single nudged-elastic-band barrier ($E_a = \SI{0.5106}{eV}$) yields a
room-temperature Cu adatom diffusion length of \SIrange{20}{30}{nm} that both
calibrates and bounds a ballistic GLAD simulator, and simultaneously specifies
why the simulator's grid-local diffusion kernel -- a bounded, non-tunable,
\SI{1.5}{nm}-reach correction -- must be replaced by a range-extended kernel
before that anchor can be applied in a growth run.

% ============================================================
\section*{Declaration on the use of AI tools}
% DRAFT 2026-09-10 -- authors to confirm model list, task scope, placement before
% submission. Policy basis: ../guide of writing articles with AI/WRITING_AND_AI_DISCLOSURE_GUIDE.md sec.8.
During the preparation of this work the authors used large language models
(Anthropic Claude; and, for bulk literature summarisation only, additional
models accessed through free-tier APIs) to assist with drafting and editing
prose, generating and refactoring analysis and simulation code, structuring
arguments, and summarising background literature. All AI-assisted text, code,
numerical results, and references were subsequently reviewed and verified by the
authors against primary sources or independent computation, and edited as
needed. The authors take full responsibility for the content of the publication.
No AI system is listed as an author, in line with COPE, ICMJE, and the relevant
publisher policies.

\section*{Declarations}

\begin{itemize}
\item \textbf{Funding:} No specific financial support statement is declared
  for this manuscript because no verified project number, contract number, or
  funding-agency wording was found in the local project records.
\item \textbf{Conflict of interest:} The authors declare no competing interests.
\item \textbf{Ethics approval and consent to participate:} Not applicable.
\item \textbf{Consent for publication:} Not applicable.
\item \textbf{Data availability:} The NEB input decks, the diffusion-length
  calculation, and the kernel-probe run configurations supporting the
  findings of this study are available from the corresponding author upon
  reasonable request.
\item \textbf{Materials availability:} Not applicable.
\item \textbf{Code availability:} The NEB, diffusion-length, and
  geometry--kinetics decomposition analysis code supporting the findings of
  this study are available from the corresponding author upon reasonable
  request.
\item \textbf{Author contribution:} O.B.H.A.: Conceptualization, Methodology,
  Software, Investigation, Formal analysis, Writing -- original draft. B.G.:
  Supervision, Writing -- review \& editing. F.C.A.: Supervision (principal),
  Writing -- review \& editing.
\end{itemize}

\bibliographystyle{aipnum4-1}
\bibliography{P4_references}

\end{document}
